% Part: turing-machines
% Chapter: machines-computations
% Section: unary-numbers

\documentclass[../../../include/open-logic-section]{subfiles}

\begin{document}

\olfileid{tur}{mac}{una}
\olsection{Representasi Wilangan kanthi Siji Simbol}

\begin{explain}
Mesin Turing makarya nganggo rerangken simbol kang ditulis
ing pitane. Gumantung marang alfabet kang digunakake mesin
Turing, rerangken simbol iki bisa nggambarake maneka input
lan output. Kang mesthi wigati yaiku mesin Turing kang
ngitung fungsi \emph{aritmetika}, yaiku fungsi ing wilangan
asli. Cara prasaja kanggo nggambarake wilangan wutuh positif
yaiku ngode nganggo rerangken siji jinis simbol~$\TMstroke$.
Yen $n \in \Nat$, tetepna $\TMstroke^n$ minangka rerangken
kosong yen $n = 0$, lan yen ora, minangka rerangken kang
dumadi saka persis $n$ simbol $\TMstroke$.
\end{explain}

\begin{defn}[Komputasi]
Mesin Turing~$M$ \emph{ngitung} fungsi
$f\colon \Nat^k \to \Nat$ yen lan mung yen
$M$~mandheg nalika nampa input
\[
\TMstroke^{n_1} \TMblank \TMstroke^{n_2} \TMblank \dots \TMblank \TMstroke^{n_k}
\]
kanthi output $\TMstroke^{f(n_1, \dots, n_k)}$.
\end{defn}

\begin{prob}
  Wenehana definisi kapan mesin Turing~$M$
  ngitung fungsi $f\colon \Nat^k \to \Nat^m$.
\end{prob}

\begin{ex}\ollabel{ex:adder}
\emph{Panambahan:}
Ayo kita gawe mesin kang ngitung fungsi $f(n,m) = n + m$.
Iki mbutuhake mesin kang diwiwiti nganggo rong blok
$\TMstroke$ dawa $n$ lan~$m$ ing pita, banjur mandheg
kanthi siji blok kang dumadi saka
$n+m$~$\TMstroke$. Rong blok input~$\TMstroke$
dipisahake dening~$\TMblank$, mula salah siji
carane yaiku nulis guratan ing kothak kang ngemot
$\TMblank$ lan mbusak~$\TMstroke$ pungkasan.
\begin{figure}\[
\begin{tikzpicture}[->,>=stealth',shorten >=1pt,auto,node distance=2.8cm,
                    semithick]
  \tikzstyle{every state}=[fill=none,draw=black,text=black]

  \node[initial,state]         (A)              {$q_0$};
  \node[state]         (B) [right of=A] {$q_1$};
  \node[state]         (C) [right of=B] {$q_2$};

  \path (A) edge node {\TMtrans{\TMblank}{\TMstroke}{\TMstay}} (B)
           edge [loop above] node {\TMtrans{\TMstroke}{\TMstroke}{\TMright}} (A)
        (B) edge [loop above] node {\TMtrans{\TMstroke}{\TMstroke}{\TMright}} (B)
            edge node {\TMtrans{\TMblank}{\TMblank}{\TMleft}} (C)
        (C) edge [loop above] node {\TMtrans{\TMstroke}{\TMblank}{\TMstay}} (C);
\end{tikzpicture}
\]\caption{Mesin kang ngitung $f(x,y) = x+y$}
\ollabel{fig:adder}
\end{figure}
\end{ex}

\begin{prob}
Telusurana konfigurasi mesin ing
\olref[tur][mac][una]{ex:adder} kanggo input~$\tuple{3,2}$.
Apa kang kedadeyan yen mesin ngitung $0+0$?
\end{prob}

\begin{explain}
Ing \olref[rep]{ex:doubler}, kita menehi tuladha
mesin Turing kang nampa rerangken~$\TMstroke$
minangka input lan mandheg kanthi rerangken
$\TMstroke$ kang cacahe tikel loro ing pita
---mesin pangganda. Nanging amarga outpute
ngemot $\TMblank$ ing sisih kiwa blok
$\TMstroke$ kang wis digandakake, mesin iku
sejatine ora ngitung fungsi $f(x) = 2x$,
senajan mbokmenawa kowe nganggep mangkono.
Kita bakal njlentrehake rong cara kanggo
mbenerake iki.
\end{explain}

\begin{ex}
Mesin ing \olref{fig:doubler-disc} ngitung fungsi
$f(x) = 2x$. Beda karo tumindake mbusak input
lan nulis rong $\TMstroke$ ing sisih tengen
adoh kanggo saben $\TMstroke$ ing input
kaya mesin saka \olref[rep]{ex:doubler},
mesin iki nambah siji~$\TMstroke$ ing sisih
tengen kanggo saben~$\TMstroke$ ing input.
Mesin kudu ngerti papan pungkasane input,
mula ditinggalake~$\TMblank$ ing antarane
input lan guratan tambahan, kang ing pungkasan
diisi~$\TMstroke$. Kita uga kudu
``ngelingi'' papan saiki ing input,
mula sawatara wektu siji~$\TMstroke$ ing
blok input diganti dadi~$\TMblank$.
  \begin{figure}
\[
\begin{tikzpicture}[->,>=stealth',shorten >=1pt,auto,node distance=2.8cm,
                    semithick]
  \tikzstyle{every state}=[fill=none,draw=black,text=black]

  \node[initial,state] (0)              {$q_0$};
  \node[state]         (1) [above of=0] {$q_1$};
  \node[state]         (2) [above of=1] {$q_2$};
  \node[state]         (3) [right of=2] {$q_3$};
  \node[state]         (4) [below of=3] {$q_4$};
  \node[state]         (5) [below of=4] {$q_5$};
  \node[state]         (6) [above right of=4] {$q_6$};
  \node[state]         (7) [below of=6] {$q_7$};
  \node[state]         (8) [below of=7] {$q_8$};

  \path (0) edge node {\TMtrans{\TMstroke}{\TMblank}{\TMright}} (1)
%    (0) edge node {\TMtrans{\TMblank}{\TMblank}{\TMstay}} (h)
    (1) edge [loop left] node {\TMtrans{\TMstroke}{\TMstroke}{\TMright}} (1)
      edge node {\TMtrans{\TMblank}{\TMblank}{\TMright}} (2)
    (2) edge [loop above] node {\TMtrans{\TMstroke}{\TMstroke}{\TMright}} (2)
        edge node {\TMtrans{\TMblank}{\TMstroke}{\TMleft}} (3)
    (3) edge [loop above] node {\TMtrans{\TMstroke}{\TMstroke}{\TMleft}} (3)
        edge node[left] {\TMtrans{\TMblank}{\TMblank}{\TMleft}} (4)
    (4) edge        node[left] {\TMtrans{\TMstroke}{\TMstroke}{\TMleft}} (5)
    (5) edge [loop below] node {\TMtrans{\TMstroke}{\TMstroke}{\TMleft}} (5)
        edge              node {\TMtrans{\TMblank}{\TMstroke}{\TMright}} (0)
    (4) edge      node[sloped] {\TMtrans{\TMblank}{\TMstroke}{\TMright}} (6)
    (6) edge              node {\TMtrans{\TMblank}{\TMstroke}{\TMright}} (7)
    (7) edge [loop right] node {\TMtrans{\TMstroke}{\TMstroke}{\TMright}} (7)
        edge              node {\TMtrans{\TMblank}{\TMblank}{\TMleft}} (8)
    (8) edge [loop right] node {\TMtrans{\TMstroke}{\TMblank}{\TMstay}} (8);
    \end{tikzpicture}
\]
\caption{Mesin kang ngitung $f(x) = 2x$}
\ollabel{fig:doubler-disc}
\end{figure}
\end{ex}

\begin{ex}\ollabel{ex:mover}
Cara kapindho kanggo ngitung $f(x) = 2x$
yaiku tetep nganggo mesin pangganda asal,
nanging nambah kahanan lan instruksi ing
pungkasan kang mindhah blok guratan kang
wis digandakake menyang sisih kiwa paling adoh
ing pita. Mesin ing \olref{fig:mover}
nindakake bagean pungkasan iki: yen diwiwiti
ing pita kang dumadi saka blok~$\TMblank$
banjur blok~$\TMstroke$ (kanthi sirah ana
ing papan sembarang ing blok~$\TMblank$),
mesin mbusak $\TMstroke$ siji-siji lan
nulis maneh ing wiwitan pita. Supaya bisa
ngerti kapan wis rampung, mesin luwih
dhisik nandhani pungkasane blok $\TMstroke$
nganggo simbol~$\TMendtape$, kang banjur
dibusak ing pungkasan. Kita miwiti nomer
kahanan saka~$q_6$, supaya bisa ditambahake
marang mesin pangganda. Kang dibutuhake
mung instruksi tambahan $\delta(q_0,
\TMblank) = \tuple{q_6,\TMblank,\TMstay}$,
yaiku panah saka~$q_0$ menyang~$q_6$
kanthi tandha~$\TMtrans{\TMblank}{\TMblank}{\TMstay}$.
(Ana sawijining masalah cilik: mesin kang
diasilake ora lumaku kanggo input~$x=0$.
Iki kita tinggalake dadi latihan.)
\begin{figure}
  \[
  \begin{tikzpicture}[->,>=stealth',shorten >=1pt,auto,node distance=2.8cm,
                      semithick]
    \tikzstyle{every state}=[fill=none,draw=black,text=black]
    \node[initial,state] (6)              {$q_6$};
    \node[state]         (7) [right of=6] {$q_7$};
    \node[state]         (8) [right of=7] {$q_8$};
    \node[state]         (9) [below of=8] {$q_9$};
    \node[state]         (10) [left of=9] {$q_{10}$};
    \node[state]         (11) [left of=10]  {$q_{11}$};
    \node[state]         (12) [below of=11]  {$q_{12}$};
    \node[state]         (13) [right of=12] {$q_{13}$};
    \node[state]         (14) [right of=13] {$q_{14}$};
    \path
    (6)  edge node {\TMtrans{\TMblank}{\TMblank}{\TMright}} (7)
    (7)  edge [loop above] node {\TMtrans{\TMblank}{\TMblank}{\TMright}} (7)
         edge node {\TMtrans{\TMstroke}{\TMstroke}{\TMright}} (8)
    (8)  edge [loop above] node {\TMtrans{\TMstroke}{\TMstroke}{\TMright}} (8)
         edge node {\TMtrans{\TMblank}{\TMendtape}{\TMleft}} (9)
    (9)  edge [loop right] node {\TMtrans{\TMstroke}{\TMstroke}{\TMleft}} (9)
         edge node {\TMtrans{\TMblank}{\TMblank}{\TMright}} (10)
    (10) edge node[above] {\TMtrans{\TMstroke}{\TMblank}{\TMleft}} (11)
    (11) edge [loop above] node {\TMtrans{\TMblank}{\TMblank}{\TMleft}} (11)
         edge node[left] {\begin{tabular}{@{}l@{}}
          \TMtrans{\TMendtape}{\TMendtape}{\TMright}\\
          \TMtrans{\TMstroke}{\TMstroke}{\TMright}
         \end{tabular}} (12)
    (12) edge node[] {\TMtrans{\TMblank}{\TMstroke}{\TMright}} (13)
    (13) edge [loop below] node {\TMtrans{\TMblank}{\TMblank}{\TMright}} (13)
         edge node[sloped] {\TMtrans{\TMstroke}{\TMblank}{\TMleft}} (11)
    (13) edge node {\TMtrans{\TMendtape}{\TMblank}{\TMstay}} (14);
  \end{tikzpicture}
  \]
  \caption{Mindhah blok $\TMstroke$ menyang kiwa}
  \ollabel{fig:mover}
\end{figure}
\end{ex}

\begin{prob}
Ing \olref[tur][mac][una]{ex:mover}, kita njlentrehake
mesin kang nyawijikake mesin pangganda saka
\olref[tur][mac][rep]{ex:doubler} lan mesin
pamindhah saka \olref[tur][mac][una]{fig:mover}.
Apa kang kedadeyan yen mesin gabungan iki
diwiwiti nganggo input~$x=0$, yaiku pita kosong?
Kepriye carane mbenerake mesin supaya ing kahanan
iki mesin mandheg kanthi output~$2x=0$?
(Kowe kudune bisa nindakake iki kanthi nambah
siji kahanan lan siji transisi.) Cathetan koreksi sumber SOL6-F046: rujukan pangganda gabungan dibenerake marang pangganda asal ing bagean Nggambarake Mesin Turing, kaya konstruksi ing conto pamindhah.
\end{prob}

\begin{prob}
\emph{Pangurangan:} Rancangna mesin Turing kang
yen diwenehi input rong rerangken guratan ora kosong
kanthi dawa $n$ lan~$m$, ing ngendi $n >
m$, ngitung fungsi $f(n,m) = n - m$.
\end{prob}

\begin{prob}
\emph{Pepadhan:} Rancangna mesin Turing kanggo
ngitung fungsi ing ngisor iki:
\[
\fn{equality}(n,m) = 
\begin{cases}
  \text{1} & \text{yen~$n = m$} \\
  \text{0} & \text{yen~$n \neq m$}
\end{cases}
\]
ing ngendi~$n$ lan~$m \in \PosInt$.
\end{prob}

\begin{prob}
Rancangna mesin Turing kanggo ngitung fungsi
$\min(x,y)$ nalika $x$ lan $y$ wilangan wutuh
positif kang digambarake ing pita minangka
rerangken $\TMstroke$ kang dipisahake dening
$\TMblank$. Kowe oleh nganggo simbol tambahan
ing alfabet mesin.

Fungsi $\min$ milih nilai paling cilik saka
argumene, mula $\min(3,5)=3$, $\min(20,16)=16$,
lan $\min(4,4)=4$, lan sateruse.
\end{prob}

\begin{defn}
  Mesin Turing~$M$ ngitung fungsi parsial $f\colon \Nat^k
  \pto \Nat$ yen lan mung yen
  \begin{enumerate}
    \item $M$ mandheg nalika nampa input $\TMstroke^{n_1}\concat\TMblank\concat
    \dots \concat\TMblank\concat\TMstroke^{n_k}$ kanthi output $\TMstroke^{m}$ yen $f(n_1, \dots, n_k) = m$.
    \item $M$ ora mandheg babar pisan, utawa mandheg kanthi
    output kang dudu siji blok~$\TMstroke$, yen
    $f(n_1, \dots, n_k)$ ora ditegesi.
  \end{enumerate}
Cathetan panyunting: mandheg tanpa output kang ditegesi miturut konvensi panandha kiwa uga kalebu ora menehi output sah.
\end{defn}

\end{document}
